\documentclass[10pt,a4paper]{amsart}
\usepackage[margin=19mm]{geometry}
\usepackage{amsmath,amssymb,mathtools}
\usepackage[scr=rsfs,cal=euler]{mathalfa}
\usepackage{xfrac}
\usepackage[unicode,hidelinks,bookmarks=false]{hyperref}
\newcommand{\ie}{i.e.\ }
\newcommand{\cf}{cf.\ }
\newcommand{\resp}{respectively\ }
\newcommand{\Subsection}{\subsection}
\providecommand{\nobreakdash}{\nobreak\hbox{-}}
\setlength{\emergencystretch}{2em}
\multlinegap0pt
\usepackage{bm}
\usepackage{bbm}
\usepackage{dsfont}
\usepackage{xspace}
\usepackage{cancel}
\usepackage{mathtools}
\usepackage{cleveref}
\usepackage{amsthm}
\makeatletter
\@ifundefined{thm}{\newtheorem{thm}{Theorem}}{}
\@ifundefined{prop}{\newtheorem{prop}[thm]{Proposition}}{}
\@ifundefined{rmk}{\newtheorem{rmk}[thm]{Remark}}{}
\makeatother
\newcommand{\Ends}[1]{\mathbf{Ends}({#1})}
\newcommand{\N}{\mathbb{N}}
\newcommand{\comp}{\operatorname{\mathbf{Comp}}}
\newcommand{\sepmax}[1]{\mathbf{Sepmax}({#1})}
\renewcommand{\mid}{ : }
\title[Paths, Ends and The Separation Problem for Infinite Graphs]{Paths, Ends and The Separation Problem for Infinite Graphs}
\author{Nicanor Carrasco-Vargas, Valentino Delle Rose,, Crist\'obal Rojas}
\date{Source: arXiv:2409.03113v2 (CC BY 4.0)}
\begin{document}
\maketitle

\noindent\emph{Source and licence.} Paths, Ends and The Separation Problem for Infinite Graphs. Version arXiv:2409.03113v2 (CC BY 4.0), distributed under the Creative Commons Attribution 4.0 International licence, as recorded in the arXiv licence metadata for this exact version and shown on the abstract page https://arxiv.org/abs/2409.03113v2. Full rights notice: licenses/arxiv-2409.03113-CC-BY-4.0.txt.\par\smallskip

\noindent\textit{Editorial scope.} Counted unit: the Thm whose source label is \texttt{thm:magic\_function} in the article (source0.tex lines 318--322, source label \texttt{thm:magic\_function}), delivered together with the article's own complete original proof of it (source0.tex lines 323--348, one \texttt{proof} environment of 26 lines). No mathematics is written, repaired or re-derived: statement and proof are the article's own text line for line. Required context retained uncounted: prop:comp-is-upper-computable (lines 243--246); comp-function-extra-obs (lines 258--263). Every definition, display and label of those lines is retained unchanged. Omitted from this package and not redistributed: 1--242, 247--257, 264--317, 349--2224. Nothing inside a retained excerpt is omitted or altered: every retained body is delivered exactly as the article's lines are written, so no omission receipt of any kind accompanies this package. Citations: the retained excerpts carry no citation command at all, so no bibliography entry of the article is transcribed and no reference is redistributed. Reference adaptation: none was needed and none was made; every reference inside the delivered unit resolves inside it, so no unresolved reference or citation is produced. Printed slips kept literal: no printed word, symbol, hypothesis or number of the counted statement or of its proof was altered. Layout and class: the article's own class is \texttt{amsart} with packages including graphicx, amsmath, amssymb, stmaryrd, setspace, nicefrac, multicol, amsthm, enumerate, amsaddr, xcolor, tikz, tcolorbox, mathrsfs; the wrapper is the lane's pinned \texttt{amsart} editorial wrapper at 10pt on A4, which restates the theorem-like environments the retained text needs and adds the article's own macro definitions that the retained text needs (\texttt{\textbackslash Ends}, \texttt{\textbackslash N}, \texttt{\textbackslash comp}, \texttt{\textbackslash mid}, \texttt{\textbackslash sepmax}), with extra wrapper packages (bm, bbm, dsfont, xspace, cancel, mathtools, cleveref). The delivered numbering is the unit's own. Assets: no retained span contains a figure environment, a graphics-inclusion command, a PDF-page inclusion, a table or a TikZ picture; no image file is copied into this package, the package declares no asset dependency and no image file is redistributed with it. Licence: version arXiv:2409.03113v2 of this article is distributed under the Creative Commons Attribution 4.0 International licence, as recorded in the arXiv licence metadata for this exact version and shown on the abstract page https://arxiv.org/abs/2409.03113v2. A full rights notice is delivered at licenses/arxiv-2409.03113-CC-BY-4.0.txt. Characters: every retained excerpt is pure ASCII, so the delivered engine, XeLaTeX with its default Latin Modern text font, reports no missing character.
\begin{prop}
\label{prop:comp-is-upper-computable}
There is a partial computable function which on input a (program for a) connected highly computable graph $G$ and a finite set of edges $E\subset E(G)$, computes a  sequence of natural numbers $(\comp_{G,n}(E))_{n\in\N}$ whose minimum equals $\comp_{G}(E)$.
\end{prop}

\begin{rmk}\label{comp-function-extra-obs}
    If $\comp_{G,n}(E)=\comp_G(E)$, then $n$ has the property that all connected components of $G\smallsetminus E$ can be determined by ``exploring'' around $E$ at distance at most $n$. 
    
    For a more precise statement, let  $u,v$ be vertices in $G\smallsetminus E$  and incident to $E$. Then they are in the same connected component of $G\smallsetminus E$ if and only if they can be connected with a path in $G\smallsetminus E$ that stays at distance at most $n$ from $E$. 
    Furthermore, $C_u$ is infinite if and only if it contains a vertex $w$ with  $d_{G\smallsetminus E}(w,u)=n$. 
\end{rmk}

\begin{thm}\label{thm:magic_function}
    Let $G$ be a connected highly computable graph with finitely many ends. Then the map $E\mapsto \comp_G(E)$ is computable. 

    In fact, there is a partial computable function which on input a (program for a) connected highly computable graph $G$ with finitely many ends, a finite set $E \subset V(G)$, the number of ends of $G$ and an arbitrary element $W \in \sepmax{G}$, outputs $\comp_{G}(E)$. 
\end{thm}

\begin{proof}
    We prove the second part of the statement, which implies the first part of the statement. Assume we receive as inputs some description of the highly computable graph $G$, a finite set of edges $W \in \sepmax{G}$, the number $k$ of ends of $G$, and a finite set $E \subset E(G)$. 
    
    In order to compute $\comp_G(E)$ we first compute  a set $U\Subset E(G)$ with the following properties: (1) $U$ contains $E$ and $W$, (2) $G[U]$ is connected, (3) for each pair of vertices  $u,v$ in $G\smallsetminus U$ that are incident to edges from $U$, and such that $u,v$ lie in the same infinite connected component of $G\smallsetminus U$, there is a path from $u$ to $v$ that is completely contained in $G[U\smallsetminus E]$, and (4) $G\smallsetminus U$ has exactly $k$ infinite connected components, and no finite connected component.
    
    Computing a set with these properties requires some care, and we need to introduce some notation. Fix a vertex $v_0$, and for each $r\in\N$ denote by $U_r$ the edge set of the induced graph $G[\{v\in V(G)\mid d_G(v,v_0)\leq r\}]$. We let $L_r$ be the set of vertices in $G\smallsetminus U_r$ that are incident to some edge from $U_r$. Observe that we can compute $U_r$ and $L_r$ from $r$ and the other input parameters, as the graph $G$ is highly computable.    
    
    We  compute $r_0$ as the smallest natural number such that both $E$ and $W$ are contained in $U_{r_0}$.  Thus $U_{r_0}\in\sepmax{G}$ and $\comp_G(E_{r_0})=k$. Since  $k$ is part of the input, we can use \Cref{prop:comp-is-upper-computable} to  compute the smallest natural number $n_0$ such that $\comp_{G,n_0}(U_{r_0})=k$. 
    

    
    It is clear that $U_{r_0+n_0}$ has properties (1) and (2) above. We prove that it also has property (3). Let $u,v\in L_{r_0+n_0}$ belong to the same infinite connected component of $G\smallsetminus U_{r_0+n_0}$. We remark that by definition of the sets $U_r$ and $L_r$, there is a path contained in $G[U_{r_0+n_0}\smallsetminus U_{r_0}]$ that joins $u$ to some element $u'\in L_{r_0}$. Similarly, there is a path contained in $G[U_{r_0+n_0}\smallsetminus U_{r_0}]$ that joins $v$ to some element $v'\in L_{r_0}$. Since $u$ and $v$ belong to the same infinite connected component of $G\smallsetminus U_{r_0}$, it follows that the same is true for the new elements $u',v'$. Finally, \Cref{comp-function-extra-obs} shows that $u'$ and $v'$ can be joined with a path contained in $G[U_{r_0+n_0}\smallsetminus U_{r_0}]$. Thus there is a path in $G[U_{r_0+n_0}\smallsetminus U_{r_0}]$ that connects $u$ and $v$. 
    
    It only remains to address property (4) above. For this purpose we let $U_f$ be the set of edges in $G$ that lie in some finite connected component of $G\smallsetminus U_{r_0+n_0}$. The set $U_f$ can be easily computed using \Cref{comp-function-extra-obs} (that is, we know that  $U_{r_0+n_0}$ belongs to $\sepmax{G}$ and thus $\comp_G(U_{r_0+n_0})=\Ends{G}$. Using \Cref{prop:comp-is-upper-computable} we can compute some $k_0$ such that $\comp_{G,k_0}(U_{r_0+n_0})=\comp_G(U_{r_0+n_0})$.  \Cref{comp-function-extra-obs} shows that finite  connected components  $G\smallsetminus U_{r_0+n_0}$ are the same as connected components of $G[U_{r_0+n_0+k_0}]$ having no vertex in $L_{r_0+n_0+k_0}$). 


    
      The set $U=U_f\cup U_{r_0+n_0}$ the properties (1), (2), (3) and (4) defined before. We will now use this set to count the number of infinite connected components of $G\smallsetminus E$. Let $V$ be the set of all vertices in $G\smallsetminus U$ that are incident to some edge from $U$, and consider a partition  $V=V_1\sqcup \dots \sqcup V_k$ where $V_i$ is the set of elements in the $i$-th connected component of $G\smallsetminus U$. Our task is to determine for which $i\ne j$, $V_i$ and $V_j$ are in the same connected component of $G\smallsetminus E$. 
    
    
    We claim that the infinite connected component of $G\smallsetminus E$ that contains $V_i$ equals the one that contains $V_j$ if and only if there is a path within $G[U]$ that visits no edge from $E$, and that joins some vertex in $V_i$ to some vertex in $V_j$. The backward implication is obvious. For the forward implication, suppose that $V_i$ and $V_j$ lie in the same connected component of $G\smallsetminus E$. Then there is a path $p$ in $G$ whose initial vertex is in $V_i$, whose final vertex is in $V_j$, and which visits no edge from $E$. 
    %Nota: la condicion (4) arriba, que $G\smallsetminus U$ no tiene componentes finitas, es para asegurar que el path sale de $U$, y entra a una componente \textit{infinita}.
    Some segments of $p$ may escape from the graph $G[U\smallsetminus E]$ through a vertex in $V$, and then enter again to $G[U\smallsetminus E]$ through another vertex in $V$. Condition (3) in the definition of $U$ ensures that we can replace these segments by segments completely contained in $G[U\smallsetminus E]$. After replacing all these segments, we end up with a path with the same initial and final vertex as $p$, but contained in $G[U\smallsetminus E]$. This shows the forward implication. 
    
    Define an equivalence relation $\sim$ in $\{1,\dots,k\}$ as follows: $i\sim j$ when there is a path in $G[U]$ that visits no edge from $E$, and that joins some vertex in $V_i$ to some vertex in $V_j$. This relation is computable because $G[U]$ is a finite graph, and we can compute the number of  equivalence classes of $\sim$. The previous paragraph shows that this number equals $\comp_G(E)$.
\end{proof} 

\end{document}
